Agamenón y Clitemnestra

Después de estos pasajes conmovedores Homero divaga, haciendo aparecer una larga lista de notables difuntas de cuyas andanzas Odiseo da cumplida cuenta en la corte de los feacios, hasta que su huésped, Alcínoo, le pregunta: 

Ah, pero dime también, y verdad verdadera te pido,			
si viste a algún divinal compañero de los que contigo
fueron un día a Ilión y encontraron allí su destino.

Ulises, por supuesto, tiene mucho que contar:


A la que acá y acullá espantó de mis ojos las ánimas			
del apretado tropel mujeril Perséfone casta,
de Agamenón el Atrida se me presentaba allí el alma
muy apenada; y con ella otras muchas allí se arrimaban
de los que en casa de Egisto hallaron con él hora mala.
Él enseguida me reconoció, que en la sangre ya andaba,			
y me lloraba en amargo quejido soltando mil lágrimas,
y a mí tendía sus manos como si tocarme anhelara;
pero ni aquel firme brío ni aquel su vigor allí estaban,
esos que antaño en sus ágiles miembros señoreaban.

Otra vez la imagen horrible del alma, que, aunque recupere brevemente el entendimiento tras beber la sangre del sacrificio, sigue siendo una sombra, sin brío ni vigor. Agamenón aparece rodeado de los hombres que le acompañaban cuando Egisto, el amante de su esposa Clitmenestra le tendió la emboscada que acabó con su vida, relatada en el canto IV por Menelao. Ulises, que no conoce el triste fin del caudillo se queda de piedra. 

Mi corazón se llenó de tristeza y al verlo lloraba,				
y le decía, alzando la voz de verbos alada:
‘Agamenón, de guerreros señor, gloriosísimo Atrida,
¿qué malhadar del morir sombriluengo domó esa tu vida?
¿Lo hizo en tus naves quizá Poseidón cuando en ellas volvías
y hórrido aliento de pérfido viento te echó encima un día?			
¿O fueron hombres, quizá, desalmados del mar a la orilla
porque sus buenos rebaños de ovejas y bueys les abrías,
o porque por su ciudad, o sus hembras quizá, combatías?’.

En fin, las maneras estándar de morir entre los guerreros que regresaban de Troya. O naufragando, o en una refriega, causada por la mala costumbre de los árgivos de robar ganado ajeno o asaltar las ciudades que se encontraban de camino, ávidos de botín. Pero el destino del Átrida ha sido muy distinto. Agamenón contesta:

‘Sangre Laercia, Odiseo milmañas, hijo del cielo,				
no me domó Poseidón en las naves a nuestro regreso
luego de alzarnos un hórrido aliento de pérfido viento,
ni fueron hombres hostiles, ya en tierra, los que me rindieron,
que Egisto fue el que dispuso mi sino y momento postrero
y me mató con mi odiosa mujer, tras llamarme a su techo,			
mientras cenaba, tal buey que uno mata en el arrendadero.

La escena que relata no puede ser más sangrienta ---a Homero el gore se le da estupendamente---: 

Morí la más triste muerte; y, conmigo, mis compañeros,
que igual que a cerdos dientirreblancos matanza les dieron
como si en casa de un rico señor poderoso con ellos
boda o escote o banquete solemne estuvieran haciendo.	

Se queja Agamenón de que una cosa es morir en combate y otra que te degüellen como ganado, durante una cena. Tres mil años más tarde, el célebre autor de <<Juego de tronos>> se limitaría a echar mano de su lectura de la Odisea para dejar al mundo patidifuso con su propia versión del asesinato de Agamenón, la célebre boda roja, sin duda inspirada en versos como estos:
		
Sé que de muchos guerreros la muerte te vino al encuentro
en fragorosa contienda o llevándote a otro en un duelo,
mas se te habría deshecho el alma muy más viendo aquello:
cómo a redor de las mesas repletas y de los calderos
yacíamos y la sangre corría empapando aquel suelo.	

Tampoco se escapa la pobre Casandra ---que ha profetizado su propia muerte---. Es Clitemnestra quién l	a mata:
	
Fue el de la hija de Príamo el frémito más lastimero,
el de Casandra, mientras Clitemnestra le hacía degüello
sobre mi cuerpo; yo alzaba los brazos, que nada pudieron
porque de hierro muriendo yo estaba; y la de ojos de perro
no esperó allí —aunque hacer ya me viera camino al infierno—		
para sellarme la boca y mis ojos cerrar con sus dedos.

Agamenón se refiere a su esposa como δὲ κυνῶπις (dé kynôpis), <<la de ojos de perro>>, usando el mismo término que Helena se aplica a sí misma en Esparta. En ambos casos, el juicio moral es menos fuerte de lo que puede parecerle a un lector moderno. El perro homérico es el animal sin vergüenza que te mira directamente a los ojos y por tanto, la traducción más exacta del epíteto, como ya comentamos anteriormente es el de <<descarada>>.  Da que pensar. Parece que a  Agamenón le moleste más la desvergüenza de la asesina ---que lo mata sin honor, degollando de paso a Casandra, a quién el rey ha hecho su amante--- que el propio crimen. Los aqueos y su extraño sentido del honor:

No, más horrible no hay nada y tampoco nada más perro
que la mujer que acomodo da a un crimen por sus adentros
como el que aquella se avino a fraguar, deshonesto y obsceno,
de darle muerte a su esposo de ley. Me esperaba yo, ingenuo,		
ser recibido con alegría por hijos y siervos
porque volvía, pero la madama de lutos y duelos
sobre ella misma vertió la vergüenza y sobre cuantas luego,
pobres mujeres, habrán de venir, aun obrando una recto’.

Odiseo se compadece, con viril solidaridad del camarada caído: 

Así me habló, y respondiéndole entonces de nuevo le digo:		
‘Ay, a la prole de Atreo el buen Zeus clarividino
mucho la debe de odiar —basta ver de esas dos los aliños—
desde el principio, que muchos por culpa de Helena caímos,
y Clitemnestra te armaba celada, estando tú ido’.

De ahí que Agamenón le aconseja que de las mujeres no hay que fiarse:

Tal dije, y él enseguida me estaba otra vez respondiendo:		
‘Por eso nunca ni con tu mujer seas tú un bendituelo
de los que de esto y aquello —de todo— le dan argumento:
cosas le cuentas algunas, mas lleva las más en secreto;


También le confirma que en asunto de mujeres, Odiseo tiene suerte. Aunque lleva un año pasándoselo estupendamente con Circe, Penélope lo espera en casa y también le espera su hijo. Agamenón, que cae poco simpático, también se lamenta de no haber podido ver al suyo y es casi el único momento en el que casi nos compadecemos de él.  

y eso que a ti no el morir te vendrá de tu esposa, Odiseo,
que es toda juicio y solo consejo de bien lleva dentro			
la joven hija de Icario, Penélope milpensamientos.
Era recién desposada, nosotros ya estábamos yendo
para la guerra y allí la dejamos con un niño al pecho,
una criatura que ya ocupará entre los hombres asiento,
feliz; y, claro, su padre por fin lo verá a su regreso,				
y correrá él a abrazarse a su padre, como es mandamiento.
La mía no permitió, mi mujer, que del hijo que tengo
llenos quedaran mis ojos; y muerte me dio antes de eso.

Y remata:

Y aún algo más te diré, y procúralo en mientes despierto:
secreto, siempre secreto, que no a descubierta, el acerco			
haz de la nave a tu patria, que yo ya en mujeres no creo.

¡Pobre! Qué desengañado está de las féminas. Se le olvidan algunos detallitos, como su afición a apropiarse de toda mujer que se cruza en su camino, empezando por Criseida, la cautiva que se empeña en conservar al principio de la Ilíada, a pesar del exuberante rescate que ofrece su padre. No solo se niega a liberar a la muchacha, sino que vocea a quién quiera oírlo que el destino de la pobre chica será envejecer sirviéndolo, literalmente, entre el telar y su cama. De paso, también compara públicamente a su esposa Clitemnestra con Criseida y afirma que es inferior a esta en cuerpo, talle, juicio y labores. Cuando los aqueos intentan hacerle razonar y Aquiles le exige que devuelva a Criseida para evitar la ira de Apolo, Agamenón se empeña en quedarse con Briseida, el botín de guerra del héroe ---que por cierto trata a su esclava con toda gentileza--- y da pie a la divina cólera de Aquiles, que casi les cuesta perder la guerra. Homero no dice que la pobre Casandra ---también cautiva de guerra--- fuera su amante, pero, visto como se las gasta el personaje, no cabe ninguna duda. Por otra parte, Egisto lo mata como un animal y Clitemnestra asesina a sangre fría a una pobre inocente. Menuda familia. 

Agamenón quiere saber de su hijo Orestes. La ávida lectora, el intenso lector de la Odisea, ya saben que este ha completado la carnicería matando a Egisto ---en el poema no se dice que mate a su madre, eso, como tantas otras extensiones de los personajes de la inmortal épica, ocurre más tarde---, pero Ulises no está al tanto de las novedades y Agamenón se queda sin la satisfacción ---que quizás Homero no le concede a propósito--- de saber que su hijo le ha vengado.

 





